\subsubsection{大規模なアーキテクチャ活動}
アーキテクチャ設計は、ライフサイクル上の一段階を指すことがある。しかし、広い意味のアーキテクティングは、それだけではない。大規模な組織やシステムでは、個別システムのアーキテクチャが、企業アーキテクチャ、製品系列アーキテクチャ、システム・オブ・システムズのアーキテクチャと関係する。

このような文脈では、次の活動が重要になる。
\begin{itemize}
  \item \keyterm{アーキテクチャ実装管理}：実装がアーキテクチャに従っているかを確認する。
  \item \keyterm{アーキテクチャ保守}：実装後もアーキテクチャを管理し、拡張する。
  \item \keyterm{アーキテクチャ管理}：組織内の複数アーキテクチャの関係を管理する。
  \item \keyterm{アーキテクチャ知識管理}：意思決定、教訓、仕様、文書などを再利用可能な資産として残す。
\end{itemize}

大規模なアーキテクチャ活動では、個別システムの最適化だけでは不十分である。組織全体で一貫性、相互運用性、再利用性、進化可能性を保つ必要がある。
